%% ma: L12-B19-0025
%% lop: 12
%% chuong: 6
%% bai: 19
%% dang: 12.19.7
%% loai: TLN
%% muc: VD
%% dap_an: "26"
%% nguon: "(NBV)-ĐỀ SỐ 16-MỖI NGÀY 1 ĐỀ THI 2025.pdf (D16, MỖI NGÀY 1 ĐỀ - Đề số 16), Phần III Câu 4"
%% kien_thuc: "Bài 19 (công thức xác suất toàn phần và công thức Bayes)"
%% trang_thai: chinh-thuc
%% ghi_chu: "Đáp án gốc 26 đúng. Bỏ ảnh minh họa ruồi giấm (ảnh trang trí, không có dữ kiện toán). Trùng ý tưởng với dạng đã duyệt 12.19.7, nhập chính thức."
\begin{ex}
Một mẫu ruồi giấm dùng để làm thí nghiệm có hai loại: ruồi mắt đỏ và ruồi mắt trắng. Biết rằng có $65\%$ là ruồi cái và $35\%$ là ruồi đực. Tỉ lệ ruồi cái mắt đỏ là $90\%$, còn tỉ lệ ruồi đực mắt đỏ là $85\%$. Chọn ngẫu nhiên một con ruồi mắt trắng. Biết rằng xác suất để đó là ruồi đực bằng $\dfrac ab$, $a,b\in\mathbb N$ và $(a,b)=1$. Tính $b-a$.
\shortans{26}
\loigiai{Gọi $D$: ``ruồi đực'', $T$: ``ruồi mắt trắng''. $P(D)=0{,}35$, $P(\overline D)=0{,}65$, $P(T\mid D)=0{,}15$, $P(T\mid\overline D)=0{,}10$.

$P(T)=0{,}35\cdot0{,}15+0{,}65\cdot0{,}10=0{,}1175$.

Theo công thức Bayes: $P(D\mid T)=\dfrac{0{,}35\cdot0{,}15}{0{,}1175}=\dfrac{21}{47}$. Vậy $a=21$, $b=47$ và $b-a=26$.}
\end{ex}
